\section{局部流形拟合}
\label{sec:mrec-local-fit}

切空间估计只给出一个平面，而流形在邻域内是弯的。要得到穿过或逼近样本的连续曲面片，还需要描述法向上怎样随切向坐标变化。本节在估计的平面上建立局部坐标，把法向分量拟合成切向坐标的函数，再分析拟合片的误差以及多片能否拼起来。拟合目标的加权、正则化等更一般的选择，以及相应的数值优化，属于\BookChapterRef{chap:optimization}{卷五第1章“优化问题与解”}开始的优化理论，本节只用普通最小二乘。

\subsection{局部图函数的拟合}
\label{sec:mrec-graph-fit}

设已选定局部中心$\bm q\in\mathbb R^D$，例如一个样本或局部质心，以及估计子空间的标准正交基$\widehat{\bm U}\in\mathbb R^{D\times m}$。把它补成$\mathbb R^D$的标准正交基$[\widehat{\bm U},\widehat{\bm V}]$，对邻域内每个观测计算切向坐标与法向分量
\[
 \bm u_j=\widehat{\bm U}^\top(\bm x_j-\bm q)\in\mathbb R^m,\qquad
 \bm w_j=\widehat{\bm V}^\top(\bm x_j-\bm q)\in\mathbb R^{D-m}.
\]
目标是找函数$g$使$\bm w_j\approx g(\bm u_j)$，由此得到连续片
\begin{equation}
 \bm\phi(\bm u)=\bm q+\widehat{\bm U}\bm u+\widehat{\bm V}g(\bm u),\qquad\bm u\in\Omega ,
 \label{eq:mrec-patch}
\end{equation}
其中$\Omega\subseteq\mathbb R^m$是切向坐标的定义域，例如半径为$\rho$的球。

\paragraph{图表示的前提}

只有当流形在这个邻域内确实是估计平面上的图时，拟合才有意义。把$M$上的点投影到估计平面，这个映射在$\bm p$处的导数是从$T_{\bm p}M$到$\widehat S$的正交投影，它的奇异值是两个子空间主夹角的余弦。只要最大主夹角小于$\pi/2$，导数就可逆，由反函数定理，$M$在$\bm p$附近可以写成$\bm q+\widehat{\bm U}\bm u+\widehat{\bm V}\widetilde g(\bm u)$，$\widetilde g$是$C^2$的。夹角越小、邻域越小，这个图表示的定义域越大、$\widetilde g$越平缓。

\paragraph{为什么要保留常数项和一次项}

真实切空间上的图函数满足$g(\bm0)=\bm0$与$Dg(\bm0)=\bm0$，但估计平面上的$\widetilde g$一般不满足：$\bm q$未必在$M$上，所以$\widetilde g(\bm0)\neq\bm0$；估计平面有倾斜，所以$D\widetilde g(\bm0)\neq\bm0$。在例~\ref{ex:mrec-one-sided}中，估计方向的斜率约为$h/R$，抛物线在新坐标下于$\bm p$处的斜率约为$-h/R$。若强行要求一次项为零，邻域边缘$|u|\approx2h$处的拟合误差就有$2h^2/R$量级，与弯曲本身一样大。

所以把估计平面的偏差纳入模型，用含常数项、一次项和二次项的多项式拟合每个法向分量：
\[
 g(\bm u)=\bm c_0+\bm L\bm u+\bigl(\tfrac12\bm u^\top\bm Q_l\bm u\bigr)_{l=1}^{D-m},
\]
其中$\bm c_0\in\mathbb R^{D-m}$，$\bm L\in\mathbb R^{(D-m)\times m}$，每个$\bm Q_l$是$m$阶对称矩阵。每个法向分量有$1+m+m(m+1)/2$个系数，$m=1$时为$3$个，$m=2$时为$6$个。

\paragraph{系数的求解与约束}

系数关于数据是线性的。把每个$\bm u_j$处全部单项式的值排成一行，得到设计矩阵$\bm\Phi\in\mathbb R^{k\times N_c}$，$N_c=1+m+m(m+1)/2$；每个法向分量的系数是以$\bm\Phi$为设计矩阵的普通最小二乘解，其构造与核对可参考\BookChapterRef{chap:linear-solution-selection}{卷二第9章“线性方程与最小二乘的选解”}。系数唯一当且仅当$\bm\Phi$满列秩。

$\bm\Phi$满列秩有直接的几何含义：$\bm\Phi\bm a=\bm0$意味着系数为$\bm a$的二次多项式在全部$\bm u_j$处为零，所以$\bm\Phi$满列秩当且仅当没有非零的二次多项式在全部切向坐标处同时为零。$m=1$时，这要求至少有三个不同的$u_j$；$m=2$时，要求至少六个点且它们不全落在同一条二次曲线上，例如全部落在一条直线上就不行。这正是第~\ref{sec:mrec-coverage}节方向分布条件在二阶拟合中的形式。

\begin{example}[圆弧上的二次拟合]
\label{ex:mrec-circle-fit}
在半径为$R$的圆的最低点取$\bm q=\bm p=(0,0)$，切线为横轴，样本取在横坐标$u=-h,0,h$处，纵坐标为$w=c,0,c$，$c=R-\sqrt{R^2-h^2}$。三点恰好确定二次多项式，解得$g(u)=(c/h^2)u^2$，常数项与一次项为零，这是对称性的结果。

二次系数$c/h^2=1/(R+\sqrt{R^2-h^2})$，当$h\to0$时趋于$1/(2R)$，与真实图函数$R-\sqrt{R^2-u^2}=u^2/(2R)+u^4/(8R^3)+\cdots$的二次项一致。在样本处拟合与真值相等，在其他位置的差别来自四次项，$|u|\leq h$时为$h^4/R^3$量级。
\end{example}

\subsection{局部片的误差与拼接}
\label{sec:mrec-patch-error}

拟合片要与真实流形比较，就要在同一个定义域上比较同一种量。在估计平面的坐标域$\Omega$上，真实流形写成$\widetilde g$，拟合片写成$g$，法向偏差是$\sup_{\bm u\in\Omega}\lVert g(\bm u)-\widetilde g(\bm u)\rVert_2$。这个量把以下几种误差合在了一起，分析时要分别检查。

\paragraph{误差的来源}

第一是切空间估计的偏差。它本身不是拟合误差，而是改变了$\widetilde g$的形状：倾斜越大，$\widetilde g$的一次项越大、图表示的定义域越小，夹角接近$\pi/2$时图表示不复存在。模型含一次项可以吸收小的倾斜，但不能挽救大的倾斜。

第二是函数截断。$\widetilde g$一般不是二次多项式。若它是$C^3$的，在半径为$\rho$的定义域上，二次Taylor多项式的余项为$\rho^3$量级；若只知道它是$C^2$的，则只能说余项比$\rho^2$小，没有确定的阶。截断误差指$\widetilde g$中任何二次多项式都消除不了的部分，Taylor余项给出它的一个上界；最小二乘所得的多项式不一定就是Taylor多项式，但它也不可能比最好的二次多项式更接近$\widetilde g$。

第三是样本分布与观测噪声。设法向数据是$\widetilde g(\bm u_j)$加上偏差$\bm\eta_j$，偏差中既有截断余项，也有观测扰动。最小二乘系数的误差不超过$\lVert\bm\eta\rVert_2/\sigma_{\min}(\bm\Phi)$，其中$\sigma_{\min}(\bm\Phi)$是设计矩阵的最小奇异值；切向坐标挤在一起或接近落在同一条二次曲线上时，$\sigma_{\min}$小，偏差被放大。有界扰动可以有系统的方向，增加样本数并不自动使它平均掉。

\paragraph{函数值误差不控制导数}

法向偏差小，不说明拟合片的切方向或曲率也准确。设真实流形在局部是$w=0$的直线，观测沿$u$方向以间距$s$交替带有$+\sigma$与$-\sigma$的法向扰动。若用插值样本的折线作拟合片，每一点的函数值误差不超过$\sigma$，斜率却到处是$\pm2\sigma/s$；$s$越小，即样本越密，斜率误差反而越大。

固定阶数的多项式不会这样振荡，但它与$\widetilde g$之差一般不是多项式，函数值上的误差界不能直接换成导数上的误差界。要从拟合片读出切方向或曲率，需要对导数单独分析误差，函数值逼近只是其中一个条件。

\paragraph{多片的并不自动是流形}

在每个样本附近各拟合一片以后，所有片的并未必是光滑流形。相邻两片在重叠区域可能不一致，留下台阶或缝隙；来自两个相近分支的片可能彼此穿过，并集出现自交，在交线附近根本不是流形。要拼成一个整体，至少需要相邻片的坐标域相互重叠、在重叠处的坐标转换光滑且一致，并用适当的方式把各片平滑地合成一张曲面。

即使拼接成功，所得曲面与真实流形同胚，也需要额外的条件和论证，这已超出本节的范围。后面的同调恢复不以局部拟合为前提：第~\ref{sec:mrec-homology-recovery}节直接由样本的球并与复形得到拓扑信息，不需要先构造曲面。
